\chapter{Dyskretne równanie Lapunowa i stabilność Schura}
\label{ch:lin23-discrete-lyapunov}
\index{równanie Steina@Równanie Steina}
\index{równanie Lapunowa!dyskretne}
\index{stabilność Schura@Stabilność Schura}
\index{promień spektralny@Promień spektralny}

\noindent\textbf{Cel rozdziału.}
W Rozdziale~\ref{ch:lin22-sylvester}
badaliśmy równanie ciągłe
\(H^*P+PH=-Q\) i układ \(x'=Hx\).
Dla układu dyskretnego
\(x_{k+1}=Ax_k\) odpowiada mu
inne równanie macierzowe:
\(A^*PA-P=-Q\).
Wykażemy, kiedy ma ono
jednoznaczne rozwiązanie,
jak wyrazić je za pomocą szeregu
i dlaczego dodatnia określoność
rozwiązania charakteryzuje stabilność.
Porównamy dyskretyzację dokładną
i metodę Eulera z dynamiką ciągłą.

W całym rozdziale
\(n\ge1\) jest liczbą całkowitą,
a \(A,Q,P\) są macierzami stopnia \(n\)
nad \(\mathbb C\), chyba że
wyraźnie zaznaczono inne ciało.
Symbol \(A^*\) oznacza
macierz sprzężoną i transponowaną.
Macierz hermitowska \(R=R^*\)
jest \emph{dodatnio określona}, jeśli
\(z^*Rz>0\) dla każdego
\(z\in\mathbb C^n\setminus\{0\}\).
Dla macierzy rzeczywistych
\(A^*=A^{\mathsf T}\).

\section{Promień spektralny i zanik potęg}
\label{sec:lin23-schur}

\begin{definicja}[Stabilność Schura]
\label{def:lin23-schur}
Niech \(A\in\operatorname{Mat}_{n\times n}(\mathbb C)\).
Przez \(\sigma(A)\) oznaczamy
zbiór jej wartości własnych.
Liczbę
\[
 \rho(A)=\max_{\lambda\in\sigma(A)}|\lambda|
\]
nazywamy \emph{promieniem spektralnym}.
Macierz \(A\) nazywamy
\emph{macierzą Schura},
czyli \emph{spektralnie stabilną
w czasie dyskretnym}, gdy
\(\rho(A)<1\).
\end{definicja}

\begin{lemat}[Zanik potęg macierzy]
\label{lem:lin23-powers}
Dla każdej macierzy
\(A\in\operatorname{Mat}_{n\times n}(\mathbb C)\)
następujące warunki są równoważne:
\begin{enumerate}[label=\textup{(\roman*)}]
 \item \(\rho(A)<1\);
 \item \(A^k\to0\) dla \(k\to\infty\);
 \item istnieją \(M>0\) i \(r\in(0,1)\)
 takie, że
 \(\|A^k\|\le Mr^k\) dla każdego
 całkowitego \(k\ge0\),
 gdzie \(\|\cdot\|\) jest normą
 operatorową na \(\mathbb C^n\).
\end{enumerate}
Jeśli \(\rho(A)<1\), to
w \textup{(iii)} można wybrać
każde \(r\) spełniające
\(\rho(A)<r<1\).
\end{lemat}
\begin{proof}
Niech \(\rho(A)<1\).
Z twierdzenia o postaci Jordana
z Rozdziału~\ref{ch:lin17-jordan}
otrzymujemy \(A=SJS^{-1}\),
gdzie \(S\) jest odwracalna,
a \(J\) stanowi sumę bloków
\(J_d(\lambda)=\lambda I_d+N\)
z \(N^d=0\).
Ponieważ \(\lambda I_d\) i \(N\)
przemieniają się, wzór dwumianowy
daje dla każdego \(k\ge0\)
\[
 J_d(\lambda)^k
 =\sum_{j=0}^{\min(k,d-1)}
 {k\choose j}\lambda^{k-j}N^j.
\]
Wszystkie wykładniki
\(\lambda^{k-j}\) są nieujemne,
więc formuła obejmuje również
przypadek \(\lambda=0\).

Wybierzmy liczby rzeczywiste
\(\rho(A)<\delta<r<1\)
z \(\delta>0\).
Ponieważ \(j\le d-1\),
dla każdego składnika sumy mamy
\[
 {k\choose j}|\lambda|^{k-j}
 \le(1+k^{d-1})\delta^{k-j}
 \le\delta^{-(d-1)}
       (1+k^{d-1})\delta^k.
\]
Ciąg
\((1+k^{d-1})(\delta/r)^k\)
jest ograniczony: iloraz kolejnych
wyrazów dąży do \(\delta/r<1\),
więc ciąg jest ostatecznie malejący.
Stąd istnieje \(M_d>0\),
dla którego
\(\|J_d(\lambda)^k\|\le M_dr^k\).
Bloków jest skończenie wiele;
dla \(A^k=SJ^kS^{-1}\)
otrzymujemy po powiększeniu
stałej \(M\) oszacowanie (iii).

Implikacja (iii) \(\Rightarrow\) (ii)
wynika z \(r^k\to0\).
Jeśli natomiast \(A^k\to0\)
i \(Av=\lambda v\) dla \(v\ne0\),
to \(\lambda^kv=A^kv\to0\).
Wobec tego \(|\lambda|<1\).
Dla wszystkich wartości własnych
otrzymujemy (i).
\end{proof}

\begin{przyklad}[Niediagonalizowalny blok stabilny]
\label{prz:lin23-jordan}
Dla
\[
 A=\begin{pmatrix}a&1\\0&a\end{pmatrix},
 \qquad a\in\mathbb C,
\]
zapis \(A=aI+N\), \(N^2=0\),
daje dla \(k\ge1\)
\[
 A^k=\begin{pmatrix}
 a^k&ka^{k-1}\\0&a^k
 \end{pmatrix}.
\]
Gdy \(a=0\), element nadprzekątny
jest równy \(1\) dla \(k=1\)
oraz \(0\) dla \(k\ge2\).
Jeżeli \(|a|<1\), oba niezerowe
typy elementów dążą do zera.
Macierz jest więc stabilna,
mimo że nie jest diagonalizowalna.
Dla \(|a|\ge1\)
wektor \((1,0)^{\mathsf T}\)
nie jest asymptotycznie tłumiony.
\end{przyklad}

\section{Równanie Steina i jego rozwiązanie}
\label{sec:lin23-stein}

\begin{definicja}[Dyskretne równanie Lapunowa]
\label{def:lin23-stein}
Dla macierzy
\(A,Q\in\operatorname{Mat}_{n\times n}(\mathbb C)\)
równanie
\[
 A^*PA-P=-Q
\]
o niewiadomej
\(P\in\operatorname{Mat}_{n\times n}(\mathbb C)\)
nazywamy \emph{dyskretnym
równaniem Lapunowa}
albo \emph{równaniem Steina}.
Wprowadzamy endomorfizm przestrzeni
macierzy
\(\mathcal T_A(X)=A^*XA\)
oraz \(\mathcal L_A=I-\mathcal T_A\).
Równanie Steina jest równoważne
równości \(\mathcal L_A(P)=Q\).
\end{definicja}

\begin{twierdzenie}[Twierdzenie Lapunowa w czasie dyskretnym]
\label{thm:lin23-lyapunov}
Dla każdej macierzy
\(A\in\operatorname{Mat}_{n\times n}(\mathbb C)\),
gdzie \(n\ge1\),
następujące warunki są równoważne:
\begin{enumerate}[label=\textup{(\roman*)}]
 \item \(\rho(A)<1\);
 \item dla każdej hermitowskiej,
 dodatnio określonej macierzy
 \(Q\in\operatorname{Mat}_{n\times n}(\mathbb C)\)
 równanie \(A^*PA-P=-Q\)
 ma dokładnie jedno rozwiązanie
 w całej przestrzeni macierzy
 zespolonych, przy czym
 rozwiązanie jest hermitowskie
 i dodatnio określone;
 \item istnieją hermitowskie,
 dodatnio określone macierze
 \(P,Q\in\operatorname{Mat}_{n\times n}(\mathbb C)\)
 spełniające \(A^*PA-P=-Q\).
\end{enumerate}
Przy tych założeniach
dla każdego \(Q\) z \textup{(ii)}
jedyne rozwiązanie określa
bezwzględnie zbieżny szereg
\[
 P=\sum_{k=0}^{\infty}(A^*)^kQA^k.
\]
\end{twierdzenie}
\begin{proof}
\emph{Implikacja (i) \(\Rightarrow\) (ii).}
Załóżmy \(\rho(A)<1\)
i wybierzmy dowolną hermitowską
dodatnio określoną \(Q\).
Z Lematu~\ref{lem:lin23-powers}
otrzymujemy \(\|A^k\|\le Mr^k\)
dla pewnych \(M>0\), \(0<r<1\).
Ponieważ
\(\|(A^*)^k\|=\|A^k\|\),
mamy
\[
 \|(A^*)^kQA^k\|
 \le M^2\|Q\|r^{2k}.
\]
Szereg geometryczny
\(\sum_{k\ge0}r^{2k}\)
jest zbieżny, a więc
szereg macierzowy z tezy
jest bezwzględnie zbieżny.
Oznaczmy jego sumę przez \(P\).

Każdy składnik jest hermitowski,
gdyż
\(((A^*)^kQA^k)^*
 =(A^*)^kQ^*A^k\).
Zatem \(P=P^*\).
Dla każdego \(z\ne0\)
\[
 z^*Pz=\sum_{k=0}^{\infty}
       (A^kz)^*Q(A^kz)
 \ge z^*Qz>0.
\]
Nierówność jest ścisła dzięki
wyrazowi \(k=0\);
pozostałe wyrazy są nieujemne.
Dowodzi to \(P>0\),
niezależnie od odwracalności \(A\).

Zbieżność bezwzględna pozwala
pomnożyć szereg przez \(A^*\)
i \(A\) oraz przesunąć indeks:
\[
 A^*PA=\sum_{k=1}^{\infty}
       (A^*)^kQA^k=P-Q.
\]
Tak skonstruowane \(P\)
spełnia równanie Steina.

Niech \(P_1,P_2\) będą
dowolnymi rozwiązaniami,
niekoniecznie hermitowskimi.
Ich różnica \(D=P_1-P_2\)
spełnia \(A^*DA=D\).
Indukcja daje
\(D=(A^*)^kDA^k\);
przez normowanie otrzymujemy
\[
 \|D\|\le M^2r^{2k}\|D\|.
\]
Przechodząc do granicy
\(k\to\infty\), mamy
\(\|D\|=0\), czyli \(D=0\).
Stąd (ii).

Implikacja (ii) \(\Rightarrow\) (iii)
wynika z wyboru \(Q=I_n\).

\emph{Implikacja (iii) \(\Rightarrow\) (i).}
Niech \(P,Q\) będą dodatnio
określone oraz
\(A^*PA-P=-Q\).
Dla dowolnej wartości własnej
\(\lambda\in\sigma(A)\)
wybierzmy \(v\ne0\)
takie, że \(Av=\lambda v\).
Wtedy
\[
 (|\lambda|^2-1)v^*Pv
 =v^*(A^*PA-P)v
 =-v^*Qv<0.
\]
Ponieważ \(v^*Pv>0\),
wynika stąd \(|\lambda|<1\).
Dotyczy to każdej wartości własnej,
a zatem \(\rho(A)<1\).
\end{proof}

\begin{uwaga}[Singularne macierze stabilne]
\label{rem:lin23-singular}
Wyraz początkowy szeregu
równa się \(Q\), toteż
\(P\ge Q>0\) w porządku
macierzy hermitowskich.
Nie jest potrzebne założenie
\(\det A\ne0\).
Dotyczy to zwłaszcza
macierzy nilpotentnych.
\end{uwaga}

\begin{przyklad}[Macierz nilpotentna]
\label{prz:lin23-nilpotent}
Dla
\[
 A=\begin{pmatrix}0&1\\0&0\end{pmatrix},
 \qquad Q=I_2,
\]
mamy \(A^2=0\),
a więc \(\rho(A)=0\).
Szereg z
Twierdzenia~\ref{thm:lin23-lyapunov}
ma tylko dwa niezerowe składniki:
\[
 P=I_2+A^{\mathsf T}A
 =\begin{pmatrix}1&0\\0&2\end{pmatrix}.
\]
Rzeczywiście,
\(A^{\mathsf T}PA
=\operatorname{diag}(0,1)\),
więc \(A^{\mathsf T}PA-P=-I_2\).
Jest to jedyne rozwiązanie.
Dla \(x_0=(a,b)^{\mathsf T}\)
otrzymujemy \(x_1=(b,0)^{\mathsf T}\),
a następnie \(x_2=0\).
\end{przyklad}

\section{Spadek energii i tempo zbieżności}
\label{sec:lin23-energy}

\begin{twierdzenie}[Energia układu dyskretnego]
\label{thm:lin23-energy}
Niech \(A\in\operatorname{Mat}_{n\times n}(\mathbb C)\)
będzie macierzą Schura.
Niech \(P,Q\) będą hermitowskimi,
dodatnio określonymi macierzami,
dla których \(A^*PA-P=-Q\).
Dla dowolnego \(x_0\in\mathbb C^n\)
i liczb całkowitych \(k\ge0\)
zdefiniujmy
\(x_k=A^kx_0\)
oraz \(E_k=x_k^*Px_k\).
Wtedy
\[
 E_{k+1}-E_k=-x_k^*Qx_k.
\]
Jeżeli \(x_k\ne0\),
to \(E_{k+1}<E_k\).
Jeżeli \(x_k=0\), to
\(E_{k+1}=E_k=0\).
Ponadto istnieje liczba
\(\gamma\in[0,1)\),
niezależna od \(x_0\),
taka, że
\[
 0\le E_k\le\gamma^kE_0
 \qquad(k\ge0).
\]
W szczególności
\(x_k\to0\) dla \(k\to\infty\).
\end{twierdzenie}
\begin{proof}
Ponieważ \(x_{k+1}=Ax_k\),
mamy
\[
 E_{k+1}-E_k
 =x_k^*(A^*PA-P)x_k
 =-x_k^*Qx_k.
\]
Dodatnia określoność \(Q\)
daje ścisłą nierówność,
gdy \(x_k\ne0\).
Jeśli \(x_k=0\), to
także \(x_{k+1}=0\)
i obie energie znikają.

Uzasadnimy oszacowanie
geometryczne bez założenia
odwracalności \(A\).
Z twierdzenia spektralnego
dla macierzy hermitowskich
z Rozdziału~\ref{ch:lin14-spektralne}
wynikają stałe
\(q_{\min},p_{\min},p_{\max}>0\)
takie, że
\[
 q_{\min}\|z\|_2^2\le z^*Qz,
 \qquad
 p_{\min}\|z\|_2^2
 \le z^*Pz\le p_{\max}\|z\|_2^2
\]
dla każdego \(z\in\mathbb C^n\).
Stąd
\(z^*Qz\ge
(q_{\min}/p_{\max})z^*Pz\).
Połóżmy
\(\eta=q_{\min}/p_{\max}>0\).
Z równania Steina mamy
\(P-Q=A^*PA\),
a prawa strona jest
nieujemnie określona,
czyli \(z^*Qz\le z^*Pz\).
Z połączenia tych nierówności
dla dowolnego niezerowego \(z\)
wynika \(0<\eta\le1\).
Dlatego
\(\gamma=1-\eta\in[0,1)\)
i
\[
 E_{k+1}
 =E_k-x_k^*Qx_k
 \le(1-\eta)E_k=\gamma E_k.
\]
Indukcja prowadzi do
\(E_k\le\gamma^kE_0\).
Ponieważ
\(p_{\min}\|x_k\|_2^2\le E_k\),
dostajemy \(x_k\to0\).
\end{proof}

\begin{wniosek}[Suma rozpraszanej energii]
\label{cor:lin23-energy-sum}
Przy założeniach
Twierdzenia~\ref{thm:lin23-energy}
dla każdego \(x_0\in\mathbb C^n\)
mamy
\[
 \sum_{k=0}^{\infty}x_k^*Qx_k
 =x_0^*Px_0.
\]
\end{wniosek}
\begin{proof}
Sumując równości
\(E_{k+1}-E_k=-x_k^*Qx_k\)
dla \(k=0,\ldots,N\),
otrzymujemy
\[
 \sum_{k=0}^{N}x_k^*Qx_k
 =E_0-E_{N+1}.
\]
Ponieważ \(E_{N+1}\to0\)
na mocy
Twierdzenia~\ref{thm:lin23-energy},
przejście do granicy
daje tezę.
\end{proof}

\section{Jednoznaczność bez stabilności Schura}
\label{sec:lin23-invertibility}

\begin{twierdzenie}[Widmo operatora Steina]
\label{thm:lin23-stein-spectrum}
Niech
\(A\in\operatorname{Mat}_{n\times n}(\mathbb C)\)
ma wartości własne
\(\lambda_1,\ldots,\lambda_n\),
z uwzględnieniem ich krotności
algebraicznych.
Niech
\(\mathcal T_A(X)=A^*XA\)
dla wszystkich
\(X\in\operatorname{Mat}_{n\times n}(\mathbb C)\).
Wielomian charakterystyczny
tego endomorfizmu
\(n^2\)-wymiarowej przestrzeni
jest równy
\[
 \chi_{\mathcal T_A}(t)=
 \prod_{i=1}^n\prod_{j=1}^n
           (t-\overline{\lambda_i}\lambda_j).
\]
Następujące warunki są równoważne:
\begin{enumerate}[label=\textup{(\roman*)}]
 \item dla każdej macierzy \(Q\)
 równanie \(A^*PA-P=-Q\)
 ma dokładnie jedno rozwiązanie \(P\);
 \item \(I-\mathcal T_A\)
 jest odwracalnym endomorfizmem
 przestrzeni macierzy;
 \item dla wszystkich
 \(i,j\in\{1,\ldots,n\}\)
 zachodzi
 \(\overline{\lambda_i}\lambda_j\ne1\).
\end{enumerate}
Warunki te nie implikują
stabilności Schura ani
dodatniej określoności
rozwiązania dla \(Q>0\).
\end{twierdzenie}
\begin{proof}
Z twierdzenia Schura z
Rozdziału~\ref{ch:lin14-spektralne}
istnieją macierz unitarna \(U\)
i górnotrójkątna \(T\),
dla których \(A=UTU^*\).
Elementami diagonalnymi \(T\)
są wartości własne
\(\lambda_1,\ldots,\lambda_n\).
Jeżeli \(X=UYU^*\), to
\[
 A^*XA=UT^*YT\,U^*.
\]
Izomorfizm
\(X\mapsto U^*XU\)
sprzęga zatem \(\mathcal T_A\)
z operatorem
\(\mathcal T_T(Y)=T^*YT\).

Niech \(E_{ij}\)
będzie jednostką macierzową.
Z górnotrójkątności \(T\)
wynika wzór
\[
 T^*E_{ij}T
 =\sum_{p=i}^{n}\sum_{q=j}^{n}
           \overline{T_{ip}}T_{jq}E_{pq}.
\]
Aby go uzasadnić, zauważmy,
że współczynnik \((p,q)\)
po lewej stronie jest równy
\((T^*)_{pi}T_{jq}
 =\overline{T_{ip}}T_{jq}\).
Może być niezerowy tylko
wtedy, gdy \(p\ge i\) i \(q\ge j\).
Przy uporządkowaniu bazy
\((E_{11},E_{12},\ldots,E_{nn})\)
leksykograficznie względem par
\((i,j)\), macierz \(\mathcal T_T\)
jest trójkątna.
Współczynnik diagonalny
dla \(E_{ij}\) jest równy
\(\overline{T_{ii}}T_{jj}
 =\overline{\lambda_i}\lambda_j\).
Wyznacznik macierzy trójkątnej
jest iloczynem jej elementów
diagonalnych; otrzymujemy
podany wzór na wielomian
charakterystyczny.

Równanie Steina ma postać
\((I-\mathcal T_A)(P)=Q\).
Endomorfizm skończenie
wymiarowej przestrzeni jest
odwracalny dokładnie wtedy,
gdy dla każdego \(Q\)
ma dokładnie jeden przeciwobraz;
to daje (i) \(\Leftrightarrow\) (ii).
Z kolei (ii) oznacza,
że \(1\) nie jest wartością
własną \(\mathcal T_A\).
Obliczony wielomian
charakterystyczny daje
równoważność z (iii).
\end{proof}

\begin{przyklad}[Jednoznaczność bez stabilności]
\label{prz:lin23-unique-unstable}
Niech
\[
 A=\operatorname{diag}(2,1/3),
 \qquad Q=I_2.
\]
Iloczyny
\(\overline{\lambda_i}\lambda_j\)
są równe \(4\), \(2/3\),
\(2/3\) oraz \(1/9\).
Żaden nie jest równy \(1\).
Z Twierdzenia~\ref{thm:lin23-stein-spectrum}
równanie Steina ma dokładnie
jedno rozwiązanie dla każdego \(Q\).
Dla wybranego \(Q=I_2\)
rozwiązanie ma postać
\[
 P=\operatorname{diag}(-1/3,9/8),
\]
gdyż jego elementy diagonalne
spełniają
\(3p_{11}=-1\)
i \((-8/9)p_{22}=-1\);
elementy pozadiagonalne
muszą być zerowe, ponieważ
\(2/3-1\ne0\).
Macierz \(P\) nie jest
dodatnio określona,
bo \(e_1^*Pe_1=-1/3\).
\end{przyklad}

\begin{przyklad}[Rezonans i wiele rozwiązań]
\label{prz:lin23-resonance}
Niech
\(A=\operatorname{diag}(2,1/2)\)
i \(Q=I_2\).
Ponieważ \(2\cdot(1/2)=1\),
operator Steina nie jest
odwracalny. Dla dowolnej
macierzy \(P=(p_{ij})\)
równanie daje
\[
 3p_{11}=-1,\quad
 0p_{12}=0,\quad
 0p_{21}=0,\quad
 -\tfrac34p_{22}=-1.
\]
Zatem wszystkie rozwiązania
to macierze
\[
 P=\begin{pmatrix}
 -1/3&u\\v&4/3
 \end{pmatrix},\qquad u,v\in\mathbb C.
\]
W podrodzinie hermitowskiej
\(v=\overline u\).
Żadne rozwiązanie nie jest
dodatnio określone,
ponieważ \(p_{11}<0\).
\end{przyklad}
